%%%%%%%%%%%%%%%%%%%%%%%%%%%%%%%%%%%%%%%%%%
% LaTeX Template pour Octo
%
% Créé par Alonso Giraldo
%
%  Toutes les images doivent se trouver
%  dans un sous-dossier pictures
%
%%%%%%%%%%%%%%%%%%%%%%%%%%%%%%%%%%%%%%%%%%

%-------------------------------------------------------------------------------
%  DOCUMENT CONFIGURATIONS
%-------------------------------------------------------------------------------

\documentclass[11pt,a4paper,twoside]{report}
\usepackage[utf8]{inputenc}
\usepackage[T1]{fontenc}
\usepackage[english,francais]{babel}            % Gestion français et anglais
\usepackage[pdftex,colorlinks]{hyperref}        % Document interactif avec liens
\usepackage{listings}                           % Blocs de code avec \lstset
\usepackage{microtype}                          % Améliore la typographie
\usepackage[explicit]{titlesec}                 % Format des titres et chapitres
\usepackage{setspace}                           % Change le linespacing
\usepackage{graphicx}
\usepackage{fancyhdr}                           % Fancy page style
\usepackage{etoolbox}                           % Utilisée pour table matière
\usepackage{lastpage}                           % Reférences sur dernière page
\usepackage{index}
\usepackage{libertine}
\usepackage{mdframed}
\usepackage{colortbl}
\usepackage{geometry}
\usepackage{acronym}
\usepackage{manfnt}
\usepackage{setspace}
\usepackage{subfig}
\usepackage{enumitem}
\usepackage{afterpage}
\usepackage[parfill]{parskip}
\usepackage{listings}
\usepackage{array}

\usepackage[acronyms,style=super,order=word]{glossaries}
\usepackage[backend=bibtex,style=verbose-trad2]{biblatex}
\usepackage{csquotes}                           % Utilisé pas biblatex

%% Configuration des acronymes
\setacronymstyle{long-sc-short}
\loadglsentries{includes/acronyms.tex}
\loadglsentries{includes/glossary.tex}

\makenoidxglossaries
% \makeindex              % Décommenter au cas ou makenoidxglossaries ne trie
% pas les entrées, dans le cas contraire supprimer

% ===============================================
%   Variables pour ce document
% ===============================================

%% Ce fichier contient les commandes et les configurations du document
% ===============================================
%   Format chapitres
% ===============================================

\titleformat{\chapter}[frame]
{\normalfont}
{\filleft
\enspace \MakeUppercase{\chaptertitlename} \thechapter\enspace}
{3ex}
{\LARGE\bfseries\filright#1}
\titlespacing{\chapter}{0pt}{*-6}{*6}

% ===============================================
%   Footer chapitres
% ===============================================

\setlength{\fboxrule}{0.5 mm}
\setlength{\fboxsep}{1 mm}
\fancypagestyle{plain}{
\fancyhf{}
\fancyhead[C]{}
\fancyfoot[L]{Octo Technology}
\fancyfoot[R]{page \thepage{}}
\fancyfoot[C]{}
\renewcommand{\headrulewidth}{0pt}
\renewcommand{\footrulewidth}{1pt}}


% ===============================================
%   Header/Footer page classique
% ===============================================

\newcommand{\head}{%
  \lhead{\textbf{\title{} - v\version}}%
  \chead{}%{\secret \vspace*{1cm} }%
  \rhead{\today}
  \lfoot{Octo Technology}
  \cfoot{\vspace*{0.5cm}}
  \rfoot{page \thepage}

  \renewcommand{\headrulewidth}{1pt}
  \renewcommand{\footrulewidth}{1pt}
}
\head

% ===============================================
%   Commandes
% ===============================================

%%Fin de document
\newcommand{\tail}{%
  \vfill%
    \begin{center}%
      \textbf{--- FIN DU DOCUMENT ---}%
    \end{center}%
\vfill%
}

%%Page vide
\newcommand\blankpage{%
  \null
  \thispagestyle{empty}%
  \addtocounter{page}{-1}%
  \newpage}

\setlength{\parindent}{1.5em}

\graphicspath{{pictures/}}

\geometry{paper=a4paper,
          includeall,
          headheight=2cm,
          headsep=0.9cm,
          vmargin=2cm,
          hmargin=3cm,
          marginparsep=0cm,
          marginparwidth=0cm,
%          showframe=true,% affiche des cadres pour visualiser les marges
}

%%%%%%%% Propriétés pdf %%%%%%%%%%%
\hypersetup{
bookmarks=true,
unicode= yes,
pdftitle={},
pdfauthor={OCTO},
pdfsubject={},
pdfkeywords={},
colorlinks=true,
linkcolor=black,
citecolor=black,
filecolor=black,
urlcolor = black
}
%%%%%%%%%%%%%%%%%%%%%%%%%%%


%%%%    version du document, utilisée en 1ère de couverture et dans les en-têtes
\newcommand{\version}{1.0}
%%%%    titre
\renewcommand{\title}{Gestion des identités dans un SI existant}
%%%%    sous titre
\newcommand{\subtitle}{Rapport - Stage fin d'année - Alonso GIRALDO}
%%%%    référence du document
\newcommand{\reference}{S01}

%% Configuration de la bibliographie
\bibliography{includes/bibliography}


%-------------------------------------------------------------------------------
%  DEBUT DOCUMENT
%-------------------------------------------------------------------------------

\begin{document}

%% Les premières pages seront numérotées en chiffres romains
\pagenumbering{Roman}

% Page de garde
\afterpage{\blankpage}
%%%%%%%%%%%%%%%%%%%%%%%%%%%%
%
% Page de garde pour Octo
%
%%%%%%%%%%%%%%%%%%%%%%%%%%%%

\pagestyle{fancy}
\thispagestyle{empty}

\begin{spacing}{1.5}%
    \center
    \includegraphics[scale=0.2]{octo_logo.png}\\%
    {\fontsize{16}{18}\selectfont \textls[150]{OCTO TECHNOLOGY}}%
\end{spacing}%

\vspace{1 cm}
\noindent\begin{minipage}[t]{6cm}
    \begin{itemize}[label=, leftmargin=*, itemsep=0pt]
        \item \textbf{Majeure:} Systèmes d'Information et Génie Logiciel (SIGL)
        \item \textbf{Promo:} 2017
        \item \textbf{UID:} 65119
        \item \textbf{Dates:} 09/01/2017 au 07/07/2017
    \end{itemize}
\end{minipage}
\hfill
\begin{minipage}[t]{5.4cm}
    \begin{flushleft}
        \begin{spacing}{1.8}
            Paris, le \today\\
            \No \textsc{\reference}\\
            Nombre de pages du document : \pageref{LastPage}
        \end{spacing}
    \end{flushleft}
\end{minipage}

\vspace{1cm}

\begin{flushright}
    {\Huge \title\\}
    \hrulefill \\ \vspace{0.5 cm}
    {\LARGE \subtitle}
\end{flushright}
\vspace{1.5 cm}

\begin{center}
    \includegraphics[scale=0.1]{epita_logo.png}%
    \vspace{0.5 cm}

    \vspace{4.5cm}

    \begin{small}
        Document version \version{}
    \end{small} \\
    {\scriptsize 34 avenue de l'Opéra 75002 Paris - Tél 01.58.56.10.00}
\end{center}

\newpage


%%Affichage des listes et tables.
\tableofcontents
\newpage

%\listoftables      % décommenter pour avoir la liste des tables
%\listofscripts     % commenter pour ne pas avoir la liste des scripts

%%permet de revenir à la numérotation arabe
\pagenumbering{arabic}

% ===============================================
%   Début effectif du document
% ===============================================

\section*{Remerciements}

\clearpage

% Partie introductive
\chapter{Résumé}

Rapide description d'octo

Rapide description du stage

Pourquoi j'ai choisi ce stage
    - Raisons
    - Ce que j'attends

Les grandes lignes du déroulement du stage



\chapter{Introduction}

\section{Constat sur le domaine du consulting en France}

Test \autocite{mrx05}
test the \gls{epita} in \gls{cabbage}

\clearpage


\listoffigures
\clearpage

\printnoidxglossaries
\printbibliography

\clearpage
\tail
\end{document}
